\section{\label{sec:dataset}Dataset}

\subsection{Device types}

We generate training data for five families of silicon photonic building blocks, summarized in
Table~\ref{tab:devices}. Together these span the core functions of a photonic integrated circuit: straight propagation,
lateral routing (S-bends), power splitting (Y-branches), multimode interference (2$\times$2 MMIs), and evanescent
coupling (directional couplers). Port counts range from 2 to 4 across the device families, and each device type has
2--5 free geometric parameters.

\begin{table}[t]
\centering
\caption{Device types in the dataset. \emph{Params} is the number of free geometric parameters; \emph{Count} is the
number of unique geometries.}
\label{tab:devices}
\begin{tabular}{lccrl}
\hline\hline
Device & Ports & Params & Count & Key parameters \\
\hline
Straight waveguide & 2 & 2 & 250  & width, length \\
S-bend             & 2 & 3 & 1000 & width, offset, $R_\text{min}$ \\
Y-branch           & 3 & 5 & 1750 & junction, bend, height \\
2$\times$2 MMI     & 4 & 5 & 3500 & MMI width/length, taper \\
Dir.\ coupler      & 4 & 5 & 3500 & gap, coupling length \\
\hline
Total              &   &   & 10\,000 & \\
\hline\hline
\end{tabular}
\end{table}

All devices use silicon-on-insulator waveguides with an effective core index $n_\text{eff} \approx 2.4$ (from
pre-computed eigenmode tables for a 220\,nm slab) and SiO$_2$ cladding with $n_\text{clad} = 1.444$. Waveguide widths
range from 0.39 to 0.58\,$\mu$m, supporting single-mode operation in the C-band. The sample allocation across device
types is weighted by parameter-space dimensionality: the 5-parameter families (MMI, Y-branch, directional coupler)
receive more geometries than the 2--3 parameter devices.

\subsection{FDTD simulation}

Each geometry is simulated at three wavelengths $\lambda \in \{1.50, 1.55, 1.60\}\,\mu$m using the Meep FDTD
solver~\cite{Oskooi2010Meep}, producing 30\,000 total simulations. The simulation configuration is:

\begin{itemize}
    \item Resolution: 14\,pixels/$\mu$m ($\Delta x \approx 71.4$\,nm).
    \item PML: $5/7 \approx 0.714\,\mu$m absorbing boundary on all sides (10\,pixels at this resolution).
    \item Source: eigenmode source exciting the fundamental TE mode at a randomly selected input port. For multi-input
          devices (MMI, directional coupler), one input port is chosen per geometry and held fixed across all three
          wavelengths.
    \item Convergence: field energy decay threshold of $10^{-5}$.
    \item Interior crop: $336 \times 112$\,pixels ($24.0 \times 8.0\,\mu$m), excluding PML regions.
\end{itemize}

Per simulation we store the complex field $E_z$, the permittivity map $\varepsilon$, a source mask identifying the
excitation port, per-port binary masks (thin cross-sections at each port location), and reference S-parameters
extracted by Meep's built-in eigenmode decomposition. Data is written to compressed sharded NPZ files with a JSON
index for efficient parallel loading.

\subsection{Sampling and splitting}

Device parameters are sampled via Latin hypercube sampling (LHS) within the ranges in Table~\ref{tab:devices}, with
values quantized to a half-pixel grid $1/(2 \times \text{resolution}) \approx 0.036\,\mu$m to ensure that each
parameter change produces a measurably different subpixel-averaged permittivity map. LHS provides better coverage of
the parameter space than uniform random sampling, particularly in the high-dimensional (5-parameter) device families.

The dataset is split 80/10/10 into train/validation/test sets by geometry ID using a deterministic hash, ensuring that
all three wavelengths for a given geometry belong to the same split (no wavelength leakage). This prevents the model
from memorizing wavelength-interpolated fields for geometries seen during training.

\subsection{Phase anchoring}

Frequency-domain fields have an arbitrary global phase: if $E_z$ is a solution to the Helmholtz equation, so is
$E_z e^{i\phi}$ for any constant $\phi$. Without phase anchoring, the same device at the same wavelength could produce
training targets with different global phases, preventing the model from learning a consistent mapping.

We anchor the global phase using the field within the source region. A weighted average of unit phasors is computed in
the source mask, weighted by local field intensity:
\begin{equation}
    \phi_\text{ref} = \arg\!\left(\sum_{p \in \mathcal{M}_\text{src}} |E_z(p)|^2 \cdot
    \frac{E_z(p)}{|E_z(p)|}\right),
    \label{eq:phase_anchor}
\end{equation}
and the entire field is rotated by $e^{-i\phi_\text{ref}}$, setting the intensity-weighted phase in the source region
to approximately zero. When restricted to high-permittivity pixels ($\varepsilon > 3.0$, i.e.\ within the waveguide
core), this anchoring is robust to the source mask geometry. If the source mask is unavailable or too small, we fall
back to a fixed region-of-interest near the input waveguide. This two-stage strategy (mask-first, ROI fallback)
ensures consistent targets even after D4 augmentation, which may rotate the source to a different edge of the domain.

\subsection{D4 symmetry augmentation}

The dihedral group D4 consists of 8 isometries of the square: four rotations (0\textdegree, 90\textdegree,
180\textdegree, 270\textdegree) and four reflections (horizontal, vertical, and two diagonal
flips)~\cite{Cohen2016GroupEquivariant}. Since the Helmholtz equation is invariant under these transformations (both
$\varepsilon$ and $E_z$ transform covariantly), applying a D4 element to a simulation--solution pair produces a valid
new training sample. We represent each element as a pair $(\text{rot}_k, \text{flip}_h)$ where $\text{rot}_k \in
\{0,1,2,3\}$ counts 90\textdegree\ CCW rotations and $\text{flip}_h \in \{\text{false}, \text{true}\}$ indicates a
horizontal flip applied after rotation, yielding $4 \times 2 = 8$ unique transforms with no
duplicates~\cite{Brandstetter2022LieSymmetry}.

During training, one of the 8 D4 elements is sampled uniformly at random per batch element and applied jointly to the
field (real and imaginary parts), the permittivity map, the source mask, and all port masks. Phase anchoring is
performed \emph{after} augmentation so that the rotated source region defines the reference phase. For rectangular
grids ($H \neq W$), the 90\textdegree\ and 270\textdegree\ rotations swap spatial dimensions, which we handle by
padding to a square domain and cropping after transformation. This effectively multiplies the training set by up to
$8\times$ without additional FDTD cost, improving generalization and reducing overfitting to specific
orientations.
